\documentclass{beamer}
% COMSC-3013 Computer Architecture shared Beamer preamble.
% Week decks live one directory below weeks/ and input this file as:
%   % COMSC-3013 Computer Architecture shared Beamer preamble.
% Week decks live one directory below weeks/ and input this file as:
%   % COMSC-3013 Computer Architecture shared Beamer preamble.
% Week decks live one directory below weeks/ and input this file as:
%   \input{../_shared/beamer-preamble.tex}
% Keep the course self-contained. Change visual/build conventions here,
% not by forking one-off week themes.

\usetheme{Boadilla}
\usecolortheme{seahorse}
\setbeamertemplate{navigation symbols}{}
\setbeamertemplate{footline}[frame number]

\usepackage{amsmath}
\usepackage{booktabs}
\usepackage{graphicx}
\usepackage{listings}
\usepackage{pgfpages}

\lstset{
  basicstyle=\ttfamily\small,
  columns=fullflexible,
  keepspaces=true,
  showstringspaces=false,
  breaklines=true
}

% Speaker notes remain hidden in the ordinary student build. Define \shownotes
% before inputting the deck to create an instructor copy; see weeks/README.md.
\ifdefined\shownotes
  \setbeameroption{show notes on second screen=right}
\fi

\newcommand{\ArchTitleSlide}[4]{%
  % #1 week label, #2 title, #3 central question, #4 Wednesday prediction
  \begin{frame}
    \titlepage
  \end{frame}
  \begin{frame}{The machine question}
    \Large #3

    \vspace{1.2em}
    \normalsize\textbf{Prediction to carry into Wednesday:} #4
  \end{frame}
}

% Keep the course self-contained. Change visual/build conventions here,
% not by forking one-off week themes.

\usetheme{Boadilla}
\usecolortheme{seahorse}
\setbeamertemplate{navigation symbols}{}
\setbeamertemplate{footline}[frame number]

\usepackage{amsmath}
\usepackage{booktabs}
\usepackage{graphicx}
\usepackage{listings}
\usepackage{pgfpages}

\lstset{
  basicstyle=\ttfamily\small,
  columns=fullflexible,
  keepspaces=true,
  showstringspaces=false,
  breaklines=true
}

% Speaker notes remain hidden in the ordinary student build. Define \shownotes
% before inputting the deck to create an instructor copy; see weeks/README.md.
\ifdefined\shownotes
  \setbeameroption{show notes on second screen=right}
\fi

\newcommand{\ArchTitleSlide}[4]{%
  % #1 week label, #2 title, #3 central question, #4 Wednesday prediction
  \begin{frame}
    \titlepage
  \end{frame}
  \begin{frame}{The machine question}
    \Large #3

    \vspace{1.2em}
    \normalsize\textbf{Prediction to carry into Wednesday:} #4
  \end{frame}
}

% Keep the course self-contained. Change visual/build conventions here,
% not by forking one-off week themes.

\usetheme{Boadilla}
\usecolortheme{seahorse}
\setbeamertemplate{navigation symbols}{}
\setbeamertemplate{footline}[frame number]

\usepackage{amsmath}
\usepackage{booktabs}
\usepackage{graphicx}
\usepackage{listings}
\usepackage{pgfpages}

\lstset{
  basicstyle=\ttfamily\small,
  columns=fullflexible,
  keepspaces=true,
  showstringspaces=false,
  breaklines=true
}

% Speaker notes remain hidden in the ordinary student build. Define \shownotes
% before inputting the deck to create an instructor copy; see weeks/README.md.
\ifdefined\shownotes
  \setbeameroption{show notes on second screen=right}
\fi

\newcommand{\ArchTitleSlide}[4]{%
  % #1 week label, #2 title, #3 central question, #4 Wednesday prediction
  \begin{frame}
    \titlepage
  \end{frame}
  \begin{frame}{The machine question}
    \Large #3

    \vspace{1.2em}
    \normalsize\textbf{Prediction to carry into Wednesday:} #4
  \end{frame}
}

\title{Week 7 Monday}\subtitle{Crack Open the CPU}\author{COMSC-3013}\date{}
\begin{document}
\begin{frame}\titlepage\note{Start from an instruction promise, not a finished spaghetti diagram.}\end{frame}
\ArchQuestionSlide{What has to exist inside the CPU for one instruction to execute correctly?}{Choose one Week 6 instruction and predict its inputs, major path, operation, control choices, and changed state.}
\begin{frame}{Start from the promise}\begin{itemize}\item The ISA defines a state change.\item The datapath exists to realize that change.\item Add structures only when the familiar instruction needs them.\end{itemize}\end{frame}
\begin{frame}{The bounded path}\begin{itemize}\item PC and instruction fetch/decode\item register file / immediate\item ALU / optional data memory\item write-back and next PC\end{itemize}\end{frame}
\begin{frame}{Why the register file matters}\begin{itemize}\item Named architectural registers hold persistent state.\item Read ports supply operands.\item Write-back changes destination state.\end{itemize}\end{frame}
\begin{frame}{Why the ALU matters}\begin{itemize}\item Arithmetic and comparisons need functional hardware.\item Address calculation reuses arithmetic.\item Control selects the required operation.\end{itemize}\end{frame}
\begin{frame}{Muxes and control}\begin{itemize}\item One machine must support multiple instruction families.\item Muxes choose sources and result paths.\item Control makes those choices instruction-dependent.\end{itemize}\end{frame}
\begin{frame}{State versus plumbing}\begin{itemize}\item Registers, memory, PC behavior are ISA-visible.\item Muxes/control/ports are implementation details of the teaching model.\item Different implementations can keep the same ISA promise.\end{itemize}\end{frame}
\begin{frame}{A second instruction class}\begin{itemize}\item Load/store adds data-memory behavior.\item Branches change the next-PC choice.\item Trace only enough contrast to make control necessary.\end{itemize}\end{frame}
\begin{frame}{Reality is richer}\begin{itemize}\item Modern CPUs may speculate and execute out of order.\item Our model is intentionally bounded.\item Necessity and evidence matter more than diagram completeness.\end{itemize}\end{frame}
\begin{frame}{Wednesday}\begin{itemize}\item Reuse Week 6 disassembly.\item Trace one known instruction.\item Repair one wrong trace.\item Separate ISA-visible state from implementation detail.\end{itemize}\end{frame}
\begin{frame}{Closing question}\centering\Large Why is each major structure necessary for the promised state change?\end{frame}
\end{document}